\FloatBarrier
\section{Discussion}\label{sec:discussion}

\subsection{Why a negative result is informative here}\label{sec:whyinformative}

A finding that a model family fails is worth little unless three things are true of the test: that the rule was fixed before the scores were seen, that the test had the power to find the structure had it been there, and that the failure can be located. The first holds in the qualified form stated in Section~\ref{sec:protocols}: the scoring cores were coded before the first pass of each system, the later completions are disclosed, and none of them decides a verdict. The second is established by different instruments in the two systems, simulation under known truth on cod (power 97--99\% against the autonomous module, 9--11\% against the stock-flow module, so the verdict is strong against one and weak against the other) and the oracle on J-17 (the map nowcasts at 7.55\,ft, so the mechanism is capable of explaining the increment). The third is what Proposition~\ref{prop:increment} supplies: the cod deficit is a misspecification term concentrated in the 1991--1995 transitions, which no member of the class can produce at the recorded catch (Proposition~\ref{prop:repeller}), and the J-17 deficit is a substitution term, the variance of a flux mechanism whose lag-one autocorrelation is 0.08 persisted as if it were one. Located failures tell the modeller where to look next. On the aquifer it is within-year information: the gap between persistence and the oracle is next year's recharge, nothing available on 31 December forecasts it, and the operational question becomes at what point in the year a nowcast with partial-year fluxes beats the annual persistence forecast, which is a new evaluation protocol rather than a new module. On the stock it is the regime: the catch-supplied modules carry a short-memory growth term with one-step skill in a low-biomass regime on one specification, and a prospective test of that claim would fix the regime boundary from a criterion external to the scores, use the consistent catch-timing convention, apply the band at both horizons and require the result on both series.

\subsection{Protocols, bands and the AR(1)}\label{sec:protocoldiff}

The two protocols disagree about exactly one module, and the disagreement is instructive rather than embarrassing. Under the aquifer's point rule the AR(1) is retained by 0.39\,ft; under the cod rule's 5\% band and two-horizon requirement it is not, because 2.96\% is inside the band and the five-year ordering reverses. The uncertainty layer sides with the stricter rule in the sense that matters: the margin's interval covers zero in both bootstrap implementations, MAE is a tie, and the five-year loss is real. The simulation on cod says what the band buys in general, most of the rule's specificity for a small cost in power (SI S8), and the AR(1) case says what it buys in particular: protection against retaining a one-horizon, within-noise margin. We therefore report the AR(1) as retained under Protocol E, as the frozen rule requires, and as a weak-autocorrelation phenomenon with no transferable content, as the evidence requires. The complementary protocol lesson comes from M2m. The best one-step forecaster on J-17 was declined by a class clause outside the retention predicate; the clause was recorded in the frozen documents, but on the rule as written M2m passes, and its margin over persistence is the only one-year margin on the series that separates from noise. A protocol that can decline the best forecaster by a clause outside its predicate should say so in its verdict, and the verdict here does: no stock-flow module is retained, and the statement is an artefact of the clause rather than of the scores.

\subsection{Limitations}\label{sec:limits}

The test is on two systems, three series, two model classes and two horizons; it is a demonstration and a replication, not a survey. On cod, each specification uses a single retrospective vintage, which favours the models (the 1983--2015 overlap between vintages differs by 24.7\,kt on average and by 131\,kt in one year), catch enters as a point series although both assessments treat it as bounded and the collapse years are under-reported, the forecasts are point forecasts scored by RMSE and MAE rather than predictive distributions scored by proper rules, sample sizes are small on Specification A, the period bands were chosen after seeing the data, the archived runner's catch-timing offset moves the catch-supplied modules' one-step margins by 20--75\,kt without moving the verdict, and the operating characteristics were measured against truths inside the class and two specific truths outside it. On J-17, total-area recharge and pumpage mix the San Antonio and Uvalde pools, recharge is a stream-loss estimate, pumpage includes unreported domestic, livestock and federal use, next-year pumpage is partly an institutional response to head rather than an exogenous flux, the map is not a karst model and the lumped-versus-distributed question is inherited, the sample is 90 years with four short fixed windows, and the nClimDiv precipitation file is not redistributed, so the three precipitation columns of the registered panel are reproducible from the provenance URL rather than from the archive alone. Proposition~\ref{prop:increment} is exact but its application is stylised: the J-17 numbers in Section~\ref{sec:decomp} use full-sample rather than rolling parameters, treat the drain term as common to both forecasters, and require the orthogonality $\E[u_{t+1}m_t]=0$, which holds approximately for a mechanism whose residual is the next year's recharge surprise but is an assumption nonetheless. Remark~\ref{rem:crossover} requires stationarity, which J-17 satisfies approximately and cod does not; its failure on cod is reported as a failure. Neither protocol was prospectively registered, and nothing in this paper should be read as a claim that it was.

\section{Conclusions}\label{sec:conclusions}

Scored on the predictand that management uses, against last-value persistence and the training mean, under protocols whose scoring cores were fixed before the first pass and whose later completions are disclosed, no structural module of a surplus-production ladder is retained on the spawning-stock biomass of Northern cod on either of two assessment series, and no structural module of a one-pool water-balance ladder is retained on the annual-mean head of the Edwards Aquifer at J-17, where the one module the aquifer's point rule retains is a univariate AR(1) with a margin inside the noise. The cod modules fail because the collapse that carries the error lies outside the model class at the recorded catch and because training-mean catch drives the autonomous modules to the floor after it; the aquifer modules fail because annual recharge is near-white and a balance that persists it substitutes an uninformative value for the quantity that explains the increment, while the same balance given the year's fluxes nowcasts at 7.55\,ft. These are the two terms of one decomposition: a one-step increment forecaster gains on persistence only through what its mechanism fits of the realised increment and loses whatever its substituted driver fails to persist, the latter term being negative whenever the driver's lag-one autocorrelation is below one half and zero when the driver is replaced by its mean; and the training mean overtakes persistence at the horizon at which the predictand's autocorrelation falls below one half, which on J-17 is two years and on the non-stationary cod series never within five. The implication for practice is procedural and general even where the numbers are not. Claims that a structural element improves prediction should be made on the state, against persistence and the next-simpler model, at the horizons advice uses, under a rule declared in advance, accompanied by what the test could have seen, and, when they fail, located in the mechanism or in the driver so that the next model is built where the error is. Mechanistic fidelity is a hypothesis about accuracy. In two unrelated systems it was scored, and in both it failed for a reason that can be stated.
